\chapter{23}

\section{Zürich (CH), Samstag, 25.
August}

Nichts wird passieren.

Louis murmelte die Worte vor sich hin. Und sah seinen neuen Pass vor
sich. Mit dem Handrücken wischte er sich die feuchte Stirn ab.

Der Wagen hielt einige Meter vor ihnen rechts am Strassenrand. Louis
fingerte die Lesebrille aus der Seitentasche seiner Hose und setzte sie
hastig auf. Mit einem grossen Schritt rückwärts entfernte er sich von
den beiden. Der stehende Mann verwarf die Hände. Seine Worte kamen
halblaut und gepresst.

«Verdammter Sch-Scheissdreck. Was wollen die Arsch-sch-Arschlöcher?»

Erst stieg links ein Polizist aus, dann rechts eine Polizistin. Sie
kamen ihnen ruhig entgegen. Der Mann war älter, ging bestimmt auf die
sechzig zu, die Frau wirkte sehr jung. Louis fielen ihr hübsches Gesicht
und ihr breiter Gang auf. Der Polizist stützte seine Arme mit
eingesteckten Daumen in seinem Gürtel ab und blieb vor Louis stehen. Die
Frau stoppte einen halben Schritt hinter ihm.

«Guten Morgen, mein Name ist Beda Tobler,» der Polizist zeigte auf die
Frau, «und das ist meine Kollegin Selina Flury,» er blickte kurz zu
Louis, dann zum am Boden liegenden Mann. Schliesslich fixierte er den
Stehenden mit strengem Blick, «was ist diesmal los?»

«Er, er gibt mir das Geld nicht zurück, dieses Schwein.»

Louis glaubte, der Mann übergebe sich nächstens. Sein Gesicht war rot
angelaufen.

«Aha, kennen Sie,» der Polizist deutete auf Louis, «den jungen Mann da?»

«Nein, und will ich auch nicht, weiss nicht, was der hier will.»

«Selina, bitte halte die zwei in Schach. Wir kümmern uns gleich um sie,»
und der Polizist ging einen Schritt auf Louis zu.

«Können Sie sich ausweisen?»

«Aber natürlich.»

Louis staunte, wie fluffig ihm die Worte über die Lippen kamen.

Der Polizist schaute nach wie vor ernst. Er hatte seine Stimme eine Spur
tiefer gelegt.

Louis zwang sich ein Lächeln ab und zog seinen Pass aus der Seitentasche
seiner Hose.

«Hier bitte.»

Er reichte dem Mann das Dokument.

Der Polizist öffnete den Pass.

«Ciro Lorente?» fragte er und nachdem er das Foto angesehen hatte,
schaute er konzentriert in Louis’ Gesicht. Dann erneut auf das Passbild.
Dann wieder direkt zu Louis. Die Angaben schienen ihn zu überzeugen. Er
wirkte entspannter und reichte Louis den Pass zurück.

«Was haben Sie mit den beiden zu tun?»

«Gar nichts. Ich bin nach Zürich unterwegs und stand plötzlich vor
ihnen, und dann hat der,» er zeigte auf den Stehenden, «dem andern einen
massiven Fusstritt verpasst und ich dachte, ich muss was tun. Ich habe
\emph{Stopp} gerufen oder so etwas, ich weiss nicht mehr genau was. Und
dann wurde der noch wütender, erstmal, ja und dann hörte ich Sie
kommen.»

«Okay, gut,» der Polizist sah entspannter aus, vielleicht wünschte es
sich Louis auch nur, «Sie sind früh unterwegs.»

Er wartete. Worauf nur? Louis versuchte, ruhig zu atmen.

«Ja, tatsächlich,» gab er zurück und erinnerte sich, dass seine
Freundin-Überraschungsgeschichte bisher gut angekommen war. Er erzählte
sie mit gespielter Fröhlichkeit. In Schmalversion.

Dem Polizisten entglitt kurz das Gesicht. Fast wäre ihm ein Schmunzeln
entwischt. Er sah ein bisschen privater aus.

«Ich wünsche Ihnen eine gute Weiterreise.»

Der Polizist hob kurz die Hand zum Abschied und wandte sich um. Louis
richtete sich gerade und verliess die Szene mit schneidigen Schritten.

«Phu,» entfuhr es ihm mit lauter Stimme, sobald er ausser Hörweite war.
Das Zittern seiner Hände liess langsam nach.

Die Feuerprobe! Der Pass wirkte echt. Er hatte seine erste
Polizeikontrolle überstanden. Er zog einen Powerriegel aus der
Einkaufstüte und ging langsam weiter. Danke, Mascha, hörte er sich leise
murmeln.

Die Strasse lag angenehm kühl vor ihm.

Ein entgegen kommender Personenwagen, der an ihm vorbeifuhr, liess eine
kraushaarige Frau am Steuer erahnen. Und augenblicklich sah er Giorgia.
Wie sie den Wagen anhielt, ausstieg und freudig strahlend auf ihn zukam.
Louis schüttelte den Kopf und fasste sich ins Haar. Er würde verrückt
werden. Er musste etwas tun.

Vielleicht könnte es gelingen, die Geschichte auszuradieren. Wenn er
sich stark genug vorstellte, er habe sich alles nur eingebildet. Die
Zeit würde für ihn arbeiten. Jede Minute, jede Stunde, die die
Katastrophe in die Vergangenheit schob, würde er sicherer. Er musste
sich nur anstrengen. Louis spürte, wie sein Kiefer zu mahlen begann.

Er blieb abrupt stehen und rieb sich die Augen. Was bastelte er sich für
einen Blödsinn zusammen. Eine fremde Stimme musste den Lead übernommen
haben. Er schob die wirren Gedanken beiseite und schritt energisch
vorwärts.

Seine Erregung legte sich mit dem Gehen.

Das wenige Wasser der Sihl zog langsam neben ihm her. Im ersten Licht
des Tages glänzte sie an einzelnen Stellen. Bald würde er da sein.

Höchste Zeit, sich der vor sich hingeschobenen Aufgabe zu stellen. Jetzt
würde er Maman anrufen.

Er erreichte den Hauptbahnhof kurz nach sieben. Louis durchschritt die
Bahnhofhalle. Das emsige Treiben zwang ihn, gelegentlich auszuweichen.
Er staunte über den regen Samstag-Betrieb. Um diese frühe Zeit.

Von neuem überlegte er, wie er sich bei Maman melden und was er sagen
sollte. Während er seine neuen Trekkingschuhe schnürte, kreisten die
Gedanken weiter. Er hatte keine Ahnung, was sie wusste, was in Mailand
unterdessen geschehen war. Kurz und beruhigend sollte sein Anruf sein.
Er wollte entspannt und gesund klingen. Oder vielleicht sollte er sie
einfach einweihen? Aber die Abwehr kam heftig. Wenn er sie damit
verängstigte, sie nicht mehr schlafen konnte, richtete er noch weiteren
Schaden an. Sie verdiente es nicht. Und Louis überkam das Gefühl dieser
sicheren Liebe, die ihn schon immer getragen hatte. Diese wärmende
Verbundenheit. Er schüttelte den Kopf und wehrte sich gegen die Tränen.
Wie oft Maman zu ihm gestanden, ihn unterstützt hatte, all die Jahre.
Selbst wenn die Dinge ausweglos schienen. Selbst wenn er über seine
Jugendjahre ein einziger Terrorist gewesen sei, wie sie später grinsend
sagte, wenn sie über vergangene Zeiten sprachen.

Sein Gehirn begann kurz im Leeren zu drehen. Er wollte sie wenigstens
nicht anlügen. Die Strandschuhe warf er in den nächsten Abfalleimer.

Louis stand inmitten der vielen Reisenden. Er versuchte, sich zu
entspannen. Noch nie hatte er vor einem Gespräch mit seiner Mutter so
viel Anspannung empfunden. Die Menschen zogen von allen Seiten an ihm
vorbei.

\vspace{\baselineskip}
\emph{Damals, Mailand. Dreizehn.}

\emph{Maman und er sitzen nebeneinander vor dem Schreibtisch des
Schuldirektors. Louis blickt trotzig um sich. Er will nicht hier sein.
Die linke Wand ist vollgehängt. Mit Stundenplänen, Post-it-Zetteln und
vielen Adressen. Sonst ist der Raum fast kahl. Wie der Direktor. Louis
ist kalt. Maman ist nervös, lächelt gequält, als der Direktor sie
begrüsst. Louis fühlt sich schlecht. Für Maman, nicht für sich. Der Typ
interessiert ihn nicht.}

\emph{Der Direktor zählt Louis’ Vergehen auf. Die nicht oder mangelhaft
erledigten Hausaufgaben. Die unentschuldigten Absenzen. Dann wird seine
Stimme lauter, obwohl er offensichtlich versucht, sich zu beherrschen.
Aggressives Arsch, denkt Louis. Er schaut kurz zu Maman und hofft, dass
sie aushält, was kommt.}

\emph{«Die gefälschte Unterschrift hat das Fass zum Überlaufen gebracht,
Signora Vincent. Unsere Schule baut auf Vertrauen auf. Das dulden wir
nicht,» sagt er gepresst, sein Gesicht rötet sich und er ändert seine
Blickrichtung. Jetzt schaut er Louis streng an. Oder wütend. Louis ist
es einerlei. Der Direktor fährt fort.}

\emph{«Ich habe dich, Louis, mit deiner Mutter eingeladen, um zusammen…»
blablabla, denkt Louis und beginnt stumm in sich drin zu singen. Eine
rhythmische, frohe Melodie. Er wiederholt sie immer wieder. Die Worte
des Direktors erreichen ihn nicht mehr.}

\emph{Trotzdem. Louis fühlt sich zerrissen. Zwischendrin. Er schämt sich,
Maman diese Situation zuzumuten. Und erlebt sich stark, als er dem
Direktor zusieht, wie der gegen die Wand läuft.}
\vspace{\baselineskip}

Bewegungslos schwer stand Louis schliesslich vor der grossen
Anzeigetafel inmitten der Bahnhofshalle.

Die Zeilen mit den verschiedenen Abfahrtszeiten, Zielorten und
Gleisangaben rutschten mit einem Klick um eine Linie nach oben.
Zuunterst erschien eine neue.

Seine Reise würde über Bern führen. Anschliessend nach Brig ins Wallis.

Er zirkelte an den vielen Reisenden vorbei, erreichte schliesslich den
SBB-Schalter und verliess ihn mit einem Ticket in der Hand. Louis
starrte auf die Karte. Hinfahrt, einfach, eine Reise ins Ungewisse. Wann
würde er wohl wieder hierstehen und die Rückfahrt lösen? Kopfschüttelnd
murmelte er ein leises «unfassbar» vor sich hin.

Die Zugsabfahrt liess ihm noch eine knappe halbe Stunde Zeit, las er,
zog mit nervösen Fingern das Handy aus der Hosentasche und stellte sich
neben eine Kaffeebar am Rande des Areals. Maman war Frühaufsteherin. Er
hoffte inbrünstig, sie möge wach sein.
